\documentclass[12pt]{article}
\usepackage{float}
\usepackage{graphicx}
\usepackage{fontspec}
\usepackage{url}
\usepackage{apacite}
\setmainfont{Times New Roman}
\usepackage[letterpaper,top=2cm,bottom=2cm,left=2cm,right=2cm,marginparwidth=1.5cm]{geometry}

\title{Weekly assignment 6: Standard Development}
\author{Truong Khang Tobias Nguyen}
\date{}

\begin{document}

\maketitle

\section{Issues}
In standard development, I can see multiple issues that hinder their development and acceptancy, leading to organizations ultimately not following or implementing them:
\begin{itemize}
\item \textbf{Pace with emerging technologies:} Typically, the process of standards takes a long time until it is developed and rolled out for implementation. Although not a problem by itself, the problem lies within the fact that technology advances at a rapid pace nowadays while developing standards that address these technological advancements does not evolve linearly with technological advances. This becomes a problem when a standard is finally developed as it may become outdated before it can be fully implemented.
\item \textbf{Stakeholder consensus:} Developing standards involve governments, organizations, experts and practicioners. All of these stakeholders have their own interests in mind when developing standards and this can lead to conflicting interests. An example could be businesses prefering inexpensive or lax measures while experts advocate for strict controls. As every party is required to participate in the development process, disagreements delay development or creates compromises that weaken security.
\item \textbf{Third-party influence:} Following up on the previous issue, the organization responsible for developing a standard can also pursue their own agenda, which can be influenced by economic or political interests. Subsequently, the SDO may become biased and favour certain measures, technologies, vendors instead of advocating for a potential stronger solution. 
\item \textbf{Accessibility:} Implementing a standard can quickly become a huge expenditure for organizations as it requires investments in technology, training, audits and processes. Justifying such an investment can become especially difficult for small and medium-sized enterprises as they lack the resources that large organizations have and may lead to them implementing inexpensive and potentially unsafe measures instead.
\item \textbf{Skill shortage:} Some standards are highly complex, requiring skilled professionals to interpret and implement them. A shortage of such experts can lead to incorrect implementation and weak enforcement, rendering the standard ineffective.
\end{itemize}

\section{Improvements}
To improve on the standard development process and to address the aformentioned issues, the following improvements can be considered:
\begin{itemize}
\item \textbf{Iterative development:} Instead of a rigid development cycle where standards are published once, standards can be iteratively developed and updated, drawing inspiration from agile methodology. This approach allows standards to be published in small increments and adapted as new technologies emerge and thus, ensuring that they do not become outdated before they can be implemented. On another note, if standards are developed modularly, adaptions would only update individual components, rendering a complete makeover after new changes unnecessary.
\item \textbf{Stakeholder collaboration:} Improving collaboration between stakeholders is an absolute necessary to not only improve the development of standards, but it also improves the standards themselves. Establishing clear scopes, measures, responsibilites and requirements early on creates a baseline and allows conflicting interests to be identified and addressed more efficiently.
\item \textbf{Transparency and oversight:} Increasing transparency and oversight during the development process reduces the risk of bias and influence of third parties. Stakeholder must disclose relevant associations with potentially conflicting parties of interest, document the reasoning behind their decisions, and external consultation. Additionally, independent reviews by different SDOs helps to ensure that decisions are based on evidence and requirements rather than interests from third parties.
\item \textbf{Tiered approach:} To make standards more accessible to SIOs, standards can follow a tiered approach that considers an organization's size, resources, industry and risk profile. This creates a baseline of requirements that all organizations are expected meet while also allowing organizations with more resources or a higher risk profile to implement more advanced measures. 
\item \textbf{Practical guidance:} Alongside of providing the necessary requirements, standards should also provide practical guidance to help organizations interpret and implement them correctly. This can be in the form of implementation checklists, reference architectures or examples. Additionally, training material and documentation could reduce the expertise needed to correctly implement the requirements and is also a step towards consistent implementation.
\end{itemize}

\section{A "best" Approach?}
I do not believe that there is a "best" approach to develop a standard. The points I made alongside their improvements do not represent the entire standard development process but individual aspects, similar to a modular system. 
\newline
On that note, even the improvement points I made do not eliminate the issues I identified but are more of a way to alleviate them as an iterative approach introduces regular changes that may become difficult to follow, improving collaboration does not eliminate conflicting interests, transparancy does not eliminate bias, tiered approaches may promote inequality and practical guidance still requires skilled professionals.
\newline
In that sense, I argue that there is no "best" approach as the entire development process consists of multiple components that need to be addressed individually. 
\newline
That being said, I still believe that one can consider the suggestions I made for each identified issue and each improvement to the individual aspect may lead to improvements to the other aspects, improving the entire development process in itself.
\end{document} 